% =============================================================================
% Chapter: Loss Functions, Training Strategy, and Experimental Design
% For: Pixel-Wise Alignment of Infrared and Visible Imagery via Dense
%       Displacement Field Prediction
% =============================================================================

\chapter{Loss Functions, Training Strategy, and Experimental Design}
\label{ch:loss_training}

Training a neural network to predict dense per-pixel displacement fields
between infrared (IR) and visible (RGB) imagery demands careful composition
of loss functions, a principled training regimen, and a rigorous experimental
protocol.  This chapter treats each of these concerns in turn.  We denote the
predicted displacement field by $\mathbf{d} = (d_x, d_y) \in
\mathbb{R}^{H \times W \times 2}$, the ground-truth field by
$\mathbf{d}^{*}$, the predicted confidence mask by
$\hat{c} \in [0,1]^{H \times W}$, and the warping operator that resamples an
image $I$ according to $\mathbf{d}$ by $\mathcal{W}(I;\mathbf{d})$.


% ============================================================================
\section{Supervised Loss Functions for Displacement Fields}
\label{sec:supervised_loss}
% ============================================================================

When ground-truth displacement fields $\mathbf{d}^{*}$ are available---either
from manual annotation, synthetic generation, or a calibrated rig---the most
direct supervisory signal is the per-pixel deviation between predicted and
true displacements.

% ---------------------------------------------------------------------------
\subsection{L1 and L2 Regression Losses}
\label{ssec:l1l2}

The \emph{L2} (mean squared error) loss penalises the squared Euclidean norm
of the displacement residual at every pixel $\mathbf{p}$:
%
\begin{equation}
  \mathcal{L}_{\text{L2}}
  = \frac{1}{|\Omega|}\sum_{\mathbf{p}\in\Omega}
    \lVert \mathbf{d}(\mathbf{p}) - \mathbf{d}^{*}(\mathbf{p}) \rVert_{2}^{2},
  \label{eq:l2}
\end{equation}
%
where $\Omega$ is the set of valid pixels.  The L2 loss provides strong
gradients for large errors, which accelerates early convergence, but it is
sensitive to outliers: a single erroneous ground-truth label can dominate the
sum.

The \emph{L1} (mean absolute error) loss replaces the squared norm with the
unsquared norm:
%
\begin{equation}
  \mathcal{L}_{\text{L1}}
  = \frac{1}{|\Omega|}\sum_{\mathbf{p}\in\Omega}
    \lVert \mathbf{d}(\mathbf{p}) - \mathbf{d}^{*}(\mathbf{p}) \rVert_{1}.
  \label{eq:l1}
\end{equation}
%
L1 loss is more robust to outliers and yields sparser gradient updates, which
can improve convergence stability when the ground-truth field contains
annotation noise~\cite{sun2018pwcnet}.

% ---------------------------------------------------------------------------
\subsection{Robust Losses: Charbonnier and Huber}
\label{ssec:robust}

The \emph{Charbonnier} (or pseudo-Huber) penalty interpolates between L1 and
L2 behaviour:
%
\begin{equation}
  \mathcal{L}_{\text{Charb}}
  = \frac{1}{|\Omega|}\sum_{\mathbf{p}\in\Omega}
    \sqrt{\lVert \mathbf{d}(\mathbf{p}) - \mathbf{d}^{*}(\mathbf{p})\rVert^{2}
          + \epsilon^{2}},
  \label{eq:charbonnier}
\end{equation}
%
with $\epsilon$ a small constant (typically $10^{-3}$).  Near zero the loss
behaves quadratically, preserving smooth gradients; for large residuals it
transitions to a linear regime, curbing the influence of
outliers~\cite{sun2014quantitative}.  The closely related \emph{Huber} loss
achieves the same effect through a piecewise definition:
%
\begin{equation}
  \rho_{\delta}(r) =
  \begin{cases}
    \tfrac{1}{2}r^{2}, & |r| \le \delta, \\
    \delta\bigl(|r| - \tfrac{1}{2}\delta\bigr), & |r| > \delta.
  \end{cases}
  \label{eq:huber}
\end{equation}
%
In our experiments (Section~\ref{sec:experiments}), the Charbonnier loss
consistently outperforms both pure L1 and pure L2, confirming the value of
outlier-robust gradients when displacement annotations carry sub-pixel noise.

% ---------------------------------------------------------------------------
\subsection{Multi-Scale Loss}
\label{ssec:multiscale}

Modern flow architectures---FlowNet~\cite{dosovitskiy2015flownet}, PWC-Net
\cite{sun2018pwcnet}, and RAFT~\cite{teed2020raft}---produce displacement
estimates at multiple spatial resolutions.  Following the practice introduced
by FlowNet, we apply the supervised loss at each scale $s \in
\{1,\dots,S\}$ with exponentially decaying weights:
%
\begin{equation}
  \mathcal{L}_{\text{ms}}
  = \sum_{s=1}^{S} \alpha_{s}\;
    \mathcal{L}_{\text{Charb}}\!\bigl(
      \mathbf{d}^{(s)},\;\text{downsample}(\mathbf{d}^{*}, s)\bigr),
  \label{eq:multiscale}
\end{equation}
%
where $\alpha_{s} = \gamma^{S-s}$ with $\gamma \in (0,1)$ (we use
$\gamma=0.8$).  RAFT refines a single-resolution estimate through recurrent
updates; its authors propose weighting the $i$-th iteration with
$\gamma^{N-i}$, which we adopt analogously.  Multi-scale supervision
provides coarse-to-fine gradient pathways, accelerating convergence and
reducing susceptibility to large-displacement local minima.

% ---------------------------------------------------------------------------
\subsection{End-Point Error as Loss and Metric}
\label{ssec:epe}

The \emph{average end-point error} (EPE) measures the mean Euclidean distance
between predicted and ground-truth displacement vectors:
%
\begin{equation}
  \text{EPE}
  = \frac{1}{|\Omega|}\sum_{\mathbf{p}\in\Omega}
    \lVert \mathbf{d}(\mathbf{p}) - \mathbf{d}^{*}(\mathbf{p}) \rVert_{2}.
  \label{eq:epe}
\end{equation}
%
EPE is equivalent to the L1 loss applied to the Euclidean norm of the
two-component residual; it serves simultaneously as a training loss and the
primary evaluation metric in optical flow
benchmarks~\cite{butler2012sintel}.  Throughout this work, we report EPE in
pixels and use it as the checkpoint-selection criterion on the validation set.


% ============================================================================
\section{Photometric and Reconstruction Losses}
\label{sec:photometric}
% ============================================================================

In many practical settings, ground-truth displacement fields are unavailable,
or only approximate.  \emph{Photometric} losses supervise the network
indirectly: warp one modality with the predicted field and measure its
agreement with the other modality.  The fundamental challenge for IR--RGB
alignment is that pixel intensities differ dramatically across modalities, so
na\"ive L1/L2 pixel loss on raw intensities is ineffective.

% ---------------------------------------------------------------------------
\subsection{Structural Similarity (SSIM)}
\label{ssec:ssim}

The structural similarity index~\cite{wang2004ssim} compares local luminance,
contrast, and structure statistics within a sliding window.  Because it
normalises for mean and variance, it captures structural correspondence even
when the absolute intensity mapping between modalities is nonlinear:
%
\begin{equation}
  \mathcal{L}_{\text{SSIM}}
  = 1 - \text{SSIM}\!\bigl(\mathcal{W}(I_{\text{IR}};\mathbf{d}),\;
        I_{\text{RGB}}\bigr).
  \label{eq:ssim}
\end{equation}
%
While SSIM is not perfectly cross-modal, it has been shown to
outperform L1/L2 losses in multimodal registration
tasks~\cite{haskins2020deep}.  In practice, we compute SSIM on
grayscale-converted RGB images using an $11 \times 11$ Gaussian window.

% ---------------------------------------------------------------------------
\subsection{Normalised Cross-Correlation (NCC)}
\label{ssec:ncc}

Local normalised cross-correlation measures the zero-mean, unit-variance
correlation within a patch centred at each pixel.  Writing $\mu$ and $\sigma$
for the local mean and standard deviation computed over a window $W$:
%
\begin{equation}
  \text{NCC}(\mathbf{p})
  = \frac{\sum_{\mathbf{q}\in W(\mathbf{p})}
      \bigl(I_{1}(\mathbf{q})-\mu_{1}\bigr)
      \bigl(I_{2}(\mathbf{q})-\mu_{2}\bigr)}
    {\sigma_{1}\;\sigma_{2}\;|W|},
  \label{eq:ncc}
\end{equation}
%
with $I_{1} = \mathcal{W}(I_{\text{IR}};\mathbf{d})$ and $I_{2} =
I_{\text{RGB}}$.  NCC is invariant to affine intensity transformations, making
it a natural choice for cross-modal
registration~\cite{balakrishnan2019voxelmorph}.  The loss is $\mathcal{L}_{\text{NCC}} = 1 - \overline{\text{NCC}}$, averaged over all pixels.

% ---------------------------------------------------------------------------
\subsection{Mutual Information Approximations}
\label{ssec:mi}

Mutual information (MI) is the gold-standard similarity measure in classical
multimodal registration~\cite{maes1997multimodal,viola1997alignment}.  Direct
MI maximisation is non-trivial in a gradient-based framework because the joint
histogram is not naturally differentiable.  Two practical proxies exist:

\paragraph{MIND descriptor.}
The \emph{modality-independent neighbourhood descriptor}
(MIND)~\cite{heinrich2012mind} encodes each pixel as a vector of
self-similarity responses within a local neighbourhood.  Because MIND captures
internal structure rather than absolute intensity, the descriptor vectors of
corresponding IR and RGB pixels are comparable.  A pointwise L2 distance
between MIND descriptors yields a differentiable, MI-flavoured loss:
%
\begin{equation}
  \mathcal{L}_{\text{MIND}}
  = \frac{1}{|\Omega|}\sum_{\mathbf{p}\in\Omega}
    \lVert \text{MIND}(I_{1},\mathbf{p})
         - \text{MIND}(I_{2},\mathbf{p}) \rVert_{2}^{2}.
  \label{eq:mind}
\end{equation}

\paragraph{Local MI via Parzen windowing.}
Differentiable Parzen-window estimators of local MI have been proposed by
\citet{guo2019mutual} and integrated into registration networks.  While
effective, they introduce additional hyperparameters (kernel bandwidth, number
of histogram bins) and incur higher computational cost.  In our ablations, the
MIND descriptor offers a favourable trade-off between accuracy and
implementation simplicity.

% ---------------------------------------------------------------------------
\subsection{Perceptual Loss}
\label{ssec:perceptual}

Features extracted from a pretrained classification network (e.g.,
VGG-19~\cite{simonyan2015vgg}) capture mid-level structural attributes---edges,
textures, shapes---that are more modality-invariant than raw
pixels~\cite{johnson2016perceptual}.  Letting $\phi^{l}$ denote the feature
map at layer $l$, the perceptual loss is:
%
\begin{equation}
  \mathcal{L}_{\text{perc}}
  = \sum_{l \in \mathcal{S}}
    \frac{1}{C_l H_l W_l}
    \bigl\lVert \phi^{l}\!\bigl(\mathcal{W}(I_{\text{IR}};\mathbf{d})\bigr)
              - \phi^{l}(I_{\text{RGB}}) \bigr\rVert_{1},
  \label{eq:perceptual}
\end{equation}
%
where $\mathcal{S}$ is a subset of layers (we use \texttt{relu2\_2},
\texttt{relu3\_3}, and \texttt{relu4\_3}).  Mid-level layers capture
structural correspondences without being too sensitive to the low-level
intensity gap between IR and RGB~\cite{zhang2018lpips}.  When the IR input is
single-channel, we replicate it across three channels before passing to VGG.

% ---------------------------------------------------------------------------
\subsection{Gradient-Based and Edge Losses}
\label{ssec:gradient}

Edges and gradients are among the most modality-invariant visual cues: object
boundaries produce strong intensity gradients in both IR and RGB regardless of
the underlying radiometric differences.  Applying a first-order differential
operator $\nabla$ (e.g., Sobel) to both images yields:
%
\begin{equation}
  \mathcal{L}_{\text{edge}}
  = \frac{1}{|\Omega|}\sum_{\mathbf{p}\in\Omega}
    \bigl\lVert |\nabla \mathcal{W}(I_{\text{IR}};\mathbf{d})(\mathbf{p})|
              - |\nabla I_{\text{RGB}}(\mathbf{p})| \bigr\rVert_{1}.
  \label{eq:edge}
\end{equation}
%
Because the Sobel operator is a fixed-weight convolution, it is naturally
differentiable.  More sophisticated options---such as a differentiable
approximation to Canny edges or the Structured Edge Detector
(SED)~\cite{dollar2013sed}---can be substituted if higher-quality edge maps
are required.  This loss complements the perceptual loss by providing an
explicit, lightweight structural supervision signal.


% ============================================================================
\section{Regularisation Losses}
\label{sec:regularisation}
% ============================================================================

Displacement fields in natural scenes are predominantly smooth, with
discontinuities confined to object boundaries.  Regularisation losses encode
this prior directly.

% ---------------------------------------------------------------------------
\subsection{Smoothness Regularisation}
\label{ssec:smooth}

A first-order smoothness penalty discourages large spatial gradients in the
predicted displacement field:
%
\begin{equation}
  \mathcal{L}_{\text{smooth}}
  = \frac{1}{|\Omega|}\sum_{\mathbf{p}\in\Omega}
    \bigl( \lVert \nabla d_x(\mathbf{p}) \rVert_{1}
         + \lVert \nabla d_y(\mathbf{p}) \rVert_{1} \bigr).
  \label{eq:smooth}
\end{equation}
%
This is the anisotropic \emph{total variation} (TV) of the displacement
field~\cite{rudin1992tv}.  An edge-aware variant down-weights the penalty at
image edges where displacement discontinuities are expected:
$\mathcal{L}_{\text{smooth}}^{\text{ea}} = \sum_{\mathbf{p}}
  e^{-\alpha |\nabla I(\mathbf{p})|}
  \lVert \nabla \mathbf{d}(\mathbf{p}) \rVert_{1},$
following~\citet{godard2017monodepth}.  Higher-order (\emph{diffusion})
regularisers---penalising the Laplacian of $\mathbf{d}$---produce smoother
fields but can over-smooth sharp motion boundaries.

% ---------------------------------------------------------------------------
\subsection{Folding and Invertibility Regularisation}
\label{ssec:folding}

A displacement field that causes two distinct source pixels to map to the same
target location is said to \emph{fold}.  Folding manifests as a non-positive
determinant of the Jacobian of the deformation $\boldsymbol{\phi}(\mathbf{p})
= \mathbf{p} + \mathbf{d}(\mathbf{p})$:
%
\begin{equation}
  \mathcal{L}_{\text{fold}}
  = \frac{1}{|\Omega|}\sum_{\mathbf{p}\in\Omega}
    \text{ReLU}\!\bigl(-\det J_{\boldsymbol{\phi}}(\mathbf{p})\bigr),
  \label{eq:fold}
\end{equation}
%
where $J_{\boldsymbol{\phi}}$ is the $2\times 2$ Jacobian, computed via
finite differences.  This penalty is zero when the deformation is locally
diffeomorphic (i.e., $\det J > 0$ everywhere) and grows as folding
worsens~\cite{dalca2019unsupervised}.  For sub-pixel IR--RGB alignment the
displacements are typically small and folding is rare, but the regulariser
provides a useful safety net when training on challenging examples with large
misalignments.

% ---------------------------------------------------------------------------
\subsection{Cyclic Consistency}
\label{ssec:cyclic}

If the network predicts both a forward displacement $\mathbf{d}_{A\to B}$ and
a backward displacement $\mathbf{d}_{B\to A}$, their composition should
approximate the identity.  The cyclic consistency loss enforces this:
%
\begin{equation}
  \mathcal{L}_{\text{cyc}}
  = \frac{1}{|\Omega|}\sum_{\mathbf{p}\in\Omega}
    \bigl\lVert \mathbf{d}_{A\to B}(\mathbf{p})
              + \mathcal{W}(\mathbf{d}_{B\to A};\mathbf{d}_{A\to B})(\mathbf{p})
    \bigr\rVert_{1}.
  \label{eq:cyclic}
\end{equation}
%
Cyclic consistency is especially valuable in the semi-supervised or
unsupervised setting, where ground-truth displacements are
absent~\cite{meister2018unflow,jonschkowski2020matters}.  It can also be used
to generate per-pixel reliability maps: pixels where the forward--backward
residual exceeds a threshold are likely occluded or unreliable.


% ============================================================================
\section{Mask and Confidence Losses}
\label{sec:mask_loss}
% ============================================================================

The predicted confidence mask $\hat{c}$ indicates the network's belief that
the displacement estimate at each pixel is trustworthy.  Supervision of this
output can be explicit or self-supervised.

% ---------------------------------------------------------------------------
\subsection{Supervised Mask Loss}
\label{ssec:mask_supervised}

When a ground-truth mask $c^{*} \in \{0,1\}^{H\times W}$ is available (e.g.,
marking regions of sensor saturation, occlusion, or failed annotation), a
standard binary cross-entropy (BCE) loss suffices:
%
\begin{equation}
  \mathcal{L}_{\text{mask}}
  = -\frac{1}{|\Omega|}\sum_{\mathbf{p}\in\Omega}\bigl[
      c^{*}\log\hat{c} + (1-c^{*})\log(1-\hat{c})
    \bigr].
  \label{eq:bce}
\end{equation}

% ---------------------------------------------------------------------------
\subsection{Self-Supervised Confidence via Learned Attenuation}
\label{ssec:learned_attenuation}

In many cases no ground-truth confidence map exists.  A principled
self-supervised alternative, introduced by~\citet{kendall2017uncertainties}
for heteroscedastic aleatoric uncertainty, modulates the reconstruction loss
by the predicted confidence and adds a regularising log-barrier:
%
\begin{equation}
  \mathcal{L}_{\text{conf}}
  = \frac{1}{|\Omega|}\sum_{\mathbf{p}\in\Omega}\left[
      \frac{\mathcal{L}_{\text{recon}}(\mathbf{p})}{\hat{c}(\mathbf{p})}
    + \log \hat{c}(\mathbf{p})
  \right],
  \label{eq:learned_attenuation}
\end{equation}
%
where $\mathcal{L}_{\text{recon}}$ is any pixel-wise reconstruction loss
(Section~\ref{sec:photometric}).  The first term incentivises the network to
assign \emph{low} confidence (small $\hat{c}$) where the reconstruction error
is large; the second term prevents the degenerate solution of setting
$\hat{c} \to 0$ everywhere, because $\log \hat{c} \to -\infty$.  The
equilibrium teaches the network to ``explain away'' genuinely difficult
regions (sensor noise, modality-specific textures) rather than distorting the
displacement field to accommodate them.

Equivalently, if the network outputs a log-variance $s = \log\sigma^{2}$
instead of a direct confidence, Eq.~\eqref{eq:learned_attenuation} becomes
the negative log-likelihood of a Gaussian with learned variance, providing a
Bayesian motivation for the attenuation trick.


% ============================================================================
\section{Total Loss}
\label{sec:total_loss}
% ============================================================================

The composite training objective is a weighted sum of the individual terms:
%
\begin{equation}
  \mathcal{L}
  = \lambda_{\text{sup}}\,\mathcal{L}_{\text{ms}}
  + \lambda_{\text{photo}}\,\mathcal{L}_{\text{photo}}
  + \lambda_{\text{smooth}}\,\mathcal{L}_{\text{smooth}}
  + \lambda_{\text{fold}}\,\mathcal{L}_{\text{fold}}
  + \lambda_{\text{cyc}}\,\mathcal{L}_{\text{cyc}}
  + \lambda_{\text{conf}}\,\mathcal{L}_{\text{conf}},
  \label{eq:total}
\end{equation}
%
where $\mathcal{L}_{\text{photo}}$ is a combination of SSIM, NCC, MIND,
perceptual, and edge losses whose relative weights are tuned by ablation
(Section~\ref{sec:experiments}).  Loss weights
$\{\lambda_{(\cdot)}\}$ are hyperparameters; we initialise them such that each
term contributes roughly equally to the total gradient magnitude at the start
of training, then hold them fixed.  An alternative is to learn the weights
using homoscedastic task uncertainty~\cite{kendall2018multitask}, though we
find manual tuning sufficient for our setting.


% ============================================================================
\section{Training Strategy}
\label{sec:training}
% ============================================================================

% ---------------------------------------------------------------------------
\subsection{Optimiser and Learning Rate}
\label{ssec:optimiser}

We use \textbf{AdamW}~\cite{loshchilov2019adamw} ($\beta_1=0.9$,
$\beta_2=0.999$, weight decay $10^{-4}$).  The learning rate follows a
\textbf{OneCycleLR} schedule~\cite{smith2019onecycle}: a linear warm-up from
$\eta_{\min}=10^{-6}$ to $\eta_{\max}=4\times10^{-4}$ over the first 10\%
of iterations, followed by cosine annealing back to $\eta_{\min}$.  OneCycleLR
consistently outperforms both step-decay and plain cosine schedules in our
experiments, echoing findings in the optical flow
literature~\cite{teed2020raft}.

% ---------------------------------------------------------------------------
\subsection{Batch Size and Memory Considerations}
\label{ssec:batchsize}

Correlation-volume--based architectures (e.g., RAFT) construct 4D cost
volumes of size $H \times W \times H \times W$ (or sampled variants thereof).
At training resolution $480\times640$, even the all-pairs correlation in
RAFT's efficient implementation consumes several gigabytes per sample.  We
use a batch size of~6 on 40\,GB A100 GPUs, employing mixed-precision
(\texttt{torch.cuda.amp}) and gradient checkpointing to reduce peak memory.
If a larger effective batch is needed (e.g., for stable batch-normalisation
statistics), gradient accumulation over two or more micro-batches is
preferred over reducing resolution.

% ---------------------------------------------------------------------------
\subsection{Pretraining on Synthetic Optical Flow Data}
\label{ssec:pretrain}

Dense ground-truth displacement fields are expensive to obtain for real
IR--RGB pairs.  Pretraining on large-scale synthetic optical flow
datasets---\textbf{FlyingChairs}~\cite{dosovitskiy2015flownet},
\textbf{FlyingThings3D}~\cite{mayer2016things3d}, and
\textbf{Sintel}~\cite{butler2012sintel}---provides the encoder with a strong
prior for correspondence estimation.  The pretraining schedule
follows the ``C+T'' protocol of RAFT: train on FlyingChairs for 120k
iterations, then FlyingThings3D for 120k iterations.

After pretraining, we \textbf{fine-tune} on the IR--RGB dataset.  During
fine-tuning, the photometric and modality-specific losses
(Section~\ref{sec:photometric}) are activated, and the learning rate is
reduced by a factor of~5.  We observe that pretraining reduces final EPE by
20--30\% relative to training from scratch, confirming that generic
correspondence priors transfer effectively to the cross-modal setting.

% ---------------------------------------------------------------------------
\subsection{Curriculum Training}
\label{ssec:curriculum}

Real-world IR--RGB misalignments range from sub-pixel parallax to tens of
pixels of offset caused by mechanical vibration or thermal expansion.  We
adopt a \emph{curriculum} strategy~\cite{bengio2009curriculum}: during the
first third of fine-tuning, the training set is restricted to pairs with
ground-truth EPE $< 5$\,px; during the second third, the threshold rises to
15\,px; for the final third, all pairs are included.  Alternatively, when
ground-truth magnitude is unknown, we generate synthetic warp augmentations of
increasing magnitude.  Curriculum training reduces early instability and
improves final accuracy on large-displacement cases by approximately
8\% EPE.

% ---------------------------------------------------------------------------
\subsection{Data Augmentation}
\label{ssec:augmentation}

We apply the following augmentations, drawn from the RAFT training
protocol~\cite{teed2020raft}:
\begin{itemize}
  \item \textbf{Spatial:} random crops ($400\times720$), horizontal and
        vertical flips, small affine perturbations (rotation $\pm 5^{\circ}$,
        scale $0.9$--$1.1$).
  \item \textbf{Photometric (applied independently to each modality):}
        brightness/contrast jitter, Gaussian noise ($\sigma$ up to 0.04),
        and random gamma correction.
  \item \textbf{Synthetic displacement augmentation:} a thin-plate-spline
        warp with random control point displacements sampled from
        $\mathcal{N}(0, \sigma_{\text{tps}}^{2})$, applied to one image while
        the displacement field is updated accordingly.  This augmentation
        effectively increases dataset diversity and has been shown beneficial
        for flow generalisation~\cite{aleotti2021learning}.
\end{itemize}

% ---------------------------------------------------------------------------
\subsection{Validation and Checkpoint Selection}
\label{ssec:validation}

\paragraph{Sequence-level splits.}
To avoid \emph{temporal leakage}---where adjacent, near-duplicate frames
appear in both training and validation---we partition at the
\textbf{sequence level}, not the frame level.  Approximately 15\% of video
sequences (covering diverse scenes, times of day, and weather conditions)
are reserved for validation, and an additional 15\% for testing.

\paragraph{Early stopping.}
We monitor validation EPE after every epoch.  If it fails to improve for 15
consecutive epochs, training is halted.  The checkpoint with the lowest
validation EPE is selected for final evaluation, providing a principled
guard against overfitting.


% ============================================================================
\section{Experimental Design and Evaluation}
\label{sec:experiments}
% ============================================================================

% ---------------------------------------------------------------------------
\subsection{Ablation Studies}
\label{ssec:ablation}

A systematic set of ablations isolates the contribution of each design
choice:

\begin{enumerate}
  \item \textbf{Loss function comparison.}  Train with (a) L2 only, (b) L1
        only, (c) Charbonnier, (d) Charbonnier + photometric composite.
        Report EPE on the validation set.
  \item \textbf{Photometric loss variants.}  Compare SSIM, NCC, MIND,
        perceptual loss, and edge loss in isolation and in combination.
  \item \textbf{With/without confidence mask.}  Evaluate whether the learned
        attenuation mechanism (Eq.~\ref{eq:learned_attenuation}) improves
        displacement accuracy in reliable regions.
  \item \textbf{Multi-scale supervision.}  Compare single-scale versus the
        $\gamma$-weighted multi-scale loss of Eq.~\eqref{eq:multiscale}.
  \item \textbf{Regularisation terms.}  Measure the effect of smoothness,
        folding, and cyclic-consistency penalties individually and jointly.
  \item \textbf{Augmentation strategies.}  Ablate spatial augmentations,
        photometric augmentations, and synthetic displacement augmentation in
        turn.
  \item \textbf{Pretraining.}  Compare from-scratch training against the
        C+T pretraining protocol of Section~\ref{ssec:pretrain}.
\end{enumerate}

% ---------------------------------------------------------------------------
\subsection{Baseline Comparisons}
\label{ssec:baselines}

We benchmark against four representative baselines spanning classical and
learned paradigms:

\begin{description}
  \item[Global homography (OpenCV).] Detect ORB keypoints, match with
    brute-force Hamming distance, and estimate a homography via RANSAC
    (\texttt{cv2.findHomography}).  This is a zero-parameter baseline that
    captures only global affine-plus-projective misalignment.
  \item[SuperPoint + SuperGlue.] A learned feature-matching
    pipeline~\cite{detone2018superpoint,sarlin2020superglue} followed by
    thin-plate-spline interpolation to produce a dense field.  This
    represents the state of the art in sparse cross-modal matching.
  \item[RAFT (pretrained).] Off-the-shelf RAFT~\cite{teed2020raft} trained
    on FlyingThings3D, evaluated zero-shot on IR--RGB pairs.  This isolates
    the benefit of domain-specific fine-tuning.
  \item[VoxelMorph.] A U-Net--based registration
    network~\cite{balakrishnan2019voxelmorph} originally designed for medical
    image registration, retrained on our data.  This represents an
    alternative dense registration architecture.
\end{description}

% ---------------------------------------------------------------------------
\subsection{Quantitative Metrics}
\label{ssec:metrics}

\begin{itemize}
  \item \textbf{End-point error (EPE):} the primary accuracy metric
        (Eq.~\ref{eq:epe}), reported in pixels.
  \item \textbf{SSIM post-warp:} SSIM between the warped IR image and the
        RGB reference (on grayscale), measuring the structural quality of
        alignment.
  \item \textbf{PSNR on edge maps:} PSNR computed between Sobel edge
        maps of the warped IR and RGB images, providing a
        modality-invariant signal quality indicator.
  \item \textbf{Mask IoU:} Intersection-over-union between the predicted
        confidence mask (thresholded at~0.5) and a ground-truth reliability
        mask, where available.
  \item \textbf{Percentage of correct keypoints (PCK):} fraction of pixels
        with EPE below a threshold $\tau$ (we report $\tau \in
        \{1, 3, 5\}$\,px).
\end{itemize}

% ---------------------------------------------------------------------------
\subsection{Qualitative Evaluation}
\label{ssec:qualitative}

Numerical metrics alone cannot capture all alignment artefacts.  We provide
four forms of qualitative evaluation:

\begin{enumerate}
  \item \textbf{Colour overlay:} the warped IR image is shown in one
        colour channel (e.g., red) and the RGB image in the complementary
        channels (green/blue).  Misalignment appears as colour fringing at
        object boundaries.
  \item \textbf{Checkerboard blending:} alternating tiles from the warped
        IR and the RGB image, making local misalignment immediately
        perceptible at tile boundaries.
  \item \textbf{Displacement field visualisation:} the predicted
        displacement field is colour-coded using the Middlebury flow colour
        wheel~\cite{baker2011middlebury}, with magnitude mapped to saturation
        and direction to hue.
  \item \textbf{Confidence mask overlay:} the predicted confidence mask is
        displayed as a semi-transparent heat map over the aligned pair,
        highlighting regions the network deems unreliable.
\end{enumerate}

These visualisations are presented for cherry-picked successes, typical cases,
and failure modes to give the reader a comprehensive understanding of model
behaviour.


% ============================================================================
\section{Summary}
\label{sec:loss_summary}
% ============================================================================

The loss landscape for cross-modal dense displacement prediction is
necessarily composite.  Supervised losses provide the strongest direct
gradient when ground-truth fields are available, while robust variants
(Charbonnier, Huber) guard against annotation noise.  Photometric losses
bridge the modality gap through structure-aware comparisons (SSIM, NCC,
MIND, perceptual features, edge maps) and enable semi-supervised learning.
Regularisation losses embed physical priors---smoothness, diffeomorphism,
cyclic consistency---that prevent degenerate solutions.  The learned
attenuation mechanism for confidence prediction elegantly couples uncertainty
estimation with the reconstruction objective.  These components, combined
with a carefully staged training protocol (synthetic pretraining, curriculum
learning, sequence-level validation) and thorough ablation-driven
experimental design, form the methodological foundation for the alignment
system evaluated in subsequent chapters.
